\chapter{ปัญญาประดิษฐ์และการประมวลผลภาพถ่ายสมาร์ทโฟนเพื่อการประเมินสุขภาพดิน}
\label{chap:ch06}

\section*{แผนบริหารการสอนประจำบทที่ 6}
\addcontentsline{toc}{section}{แผนบริหารการสอนประจำบทที่ 6}

\noindent\textbf{หัวข้อเนื้อหาประจำบท}
\begin{enumerate}
    \item หลักการประมวลผลภาพถ่ายดิจิทัลและปริภูมิสีของดิน
    \item การเทียบค่าสีมาตรฐานด้วย Color Calibration Card
    \item สถาปัตยกรรมโมเดล Deep Learning สำหรับการทำนายคุณภาพดิน
\end{enumerate}

\noindent\textbf{วัตถุประสงค์เชิงพฤติกรรม}
\begin{enumerate}
    \item อธิบายหลักการ ทฤษฎี และสถาปัตยกรรมเทคโนโลยีประจำบทที่ 6 ได้อย่างถูกต้อง
    \item ประยุกต์ใช้เครื่องมือ อุปกรณ์เซนเซอร์ และระบบปัญญาประดิษฐ์เพื่อการจัดการสวนทุเรียนได้อย่างมีประสิทธิภาพ
    \item วิเคราะห์ สรุปผลข้อมูล และแก้ไขปัญหาเชิงฟิสิกส์เกษตรและคุณภาพดินได้อย่างเป็นระบบ
\end{enumerate}

\noindent\textbf{การวัดและประเมินผล}
\begin{table}[H]
\centering\small
\begin{tabularx}{\textwidth}{p{0.55\textwidth} p{0.25\textwidth} c}
\toprule
\textbf{เกณฑ์การประเมินผล} & \textbf{เครื่องมือวัดผล} & \textbf{สัดส่วนคะแนน} \\
\midrule
1. ทักษะการปฏิบัติงานและการใช้อุปกรณ์ & แบบประเมินทักษะภาคสนาม & 40\% \\
2. รายงานข้อมูลและการวิเคราะห์ผล & การตรวจสมุดบันทึกและดาต้าชีท & 30\% \\
3. แบบฝึกหัดและคำถามท้ายบท & แบบทดสอบท้ายบทเรียน & 20\% \\
4. การมีส่วนร่วมและความรับผิดชอบ & การเข้าเรียนและการตรงต่อเวลา & 10\% \\
\midrule
\textbf{รวมทั้งสิ้น} & & \textbf{100\%} \\
\bottomrule
\end{tabularx}
\end{table}
\newpage

% ==========================================================
% เนื้อหาบทเรียนประจำบทที่ 6
% ==========================================================

\section{หลักการประมวลผลภาพถ่ายดิจิทัลและปริภูมิสีของดิน}
\index{หลักการประมวลผลภาพถ่ายดิจิทัลและปริภูมิสีของดิน}

การศึกษาและทำความเข้าใจเกี่ยวกับหลักการประมวลผลภาพถ่ายดิจิทัลและปริภูมิสีของดิน (Digital Soil Image Processing and Color Spaces) มีความสำคัญอย่างยิ่งในระบบเกษตรอัจฉริยะ โดยมีรายละเอียดสาระสำคัญดังนี้

สีของดิน (Soil Color) จัดเป็นตัวบ่งชี้ทางกายภาพที่สำคัญซึ่งสะท้อนถึงองค์ประกอบทางเคมีและแร่ธาตุ เช่น ปริมาณอินทรียวัตถุ (ดินสีคล้ำ/ดำ), ปริมาณสารประกอบเหล็กออกไซด์ (ดินสีแดง/เหลือง), และระดับการระบายน้ำของดิน

การถ่ายภาพดินด้วยกล้องสมาร์ทโฟนทั่วไปจะบันทึกข้อมูลในปริภูมิสี RGB (Red, Green, Blue) ซึ่งมีความไวต่อการเปลี่ยนแปลงของความเข้มแสงธรรมชาติอย่างยิ่ง ดังนั้นในระบบประมวลผล SoilAI จึงต้องแปลงปริภูมิสีเข้าสู่ระบบที่มีความคงทนต่อแสง:
1. ระบบสี HSV/HSL: แยกค่า Hue (เนื้อสี), Saturation (ความอิ่มตัวสี), และ Value/Lightness (ความสว่าง)
2. ระบบสี CIE L*a*b*: ค่า L* แสดงความสว่าง (Lightness ซึ่งแปรผกผันกับปริมาณอินทรียวัตถุและความชื้นในดิน), ค่า a* แสดงความแดง-เขียว (สะท้อนแร่ฮีมาไทต์ Fe$_2$O$_3$), และค่า b* แสดงความเหลือง-น้ำเงิน (สะท้อนแร่เกอไทต์ FeOOH)
3. ระบบสี Munsell Soil Color: ระบบจำแนกสีดินมาตรฐานสากลตามค่า Hue, Value, และ Chroma

\section{การเทียบค่าสีมาตรฐานด้วย Color Calibration Card}
\index{การเทียบค่าสีมาตรฐานด้วย Color Calibration Card}

การศึกษาและทำความเข้าใจเกี่ยวกับการเทียบค่าสีมาตรฐานด้วย Color Calibration Card (Color Calibration Using Reference Cards) มีความสำคัญอย่างยิ่งในระบบเกษตรอัจฉริยะ โดยมีรายละเอียดสาระสำคัญดังนี้

เพื่อขจัดผลกระทบจากความแปรปรวนของแสงแดด เงาเมฆ และเซนเซอร์กล้องสมาร์ทโฟนแต่ละยี่ห้อ ระบบจำเป็นต้องใช้ การเทียบสีอ้างอิง (Color Calibration):
- ใช้บัตรสีมาตรฐานที่มีแถบสีเทากลาง 18\% (18\% Neutral Gray Card) และแถบแม่สี RGB วางระนาบเดียวกับผิวดิน
- อัลกอริทึม White Balance และ Color Constancy จะคำนวณเวกเตอร์แปลงค่าสี (Color Correction Matrix; CCM)
\[ \begin{bmatrix} R\_{\text{corrected}} \\ G\_{\text{corrected}} \\ B\_{\text{corrected}} \end{bmatrix} = \mathbf{M}\_{3\times 3} \begin{bmatrix} R\_{\text{raw}} \\ G\_{\text{raw}} \\ B\_{\text{raw}} \end{bmatrix} \]
ช่วยให้ภาพถ่ายดินที่ถ่ายในช่วงเวลาต่างกันหรือต่างสมาร์ทโฟนมีค่าสีดิจิทัลที่เสถียรและเทียบเคียงกันได้

\section{สถาปัตยกรรมโมเดล Deep Learning สำหรับการทำนายคุณภาพดิน}
\index{สถาปัตยกรรมโมเดล Deep Learning สำหรับการทำนายคุณภาพดิน}

การศึกษาและทำความเข้าใจเกี่ยวกับสถาปัตยกรรมโมเดล Deep Learning สำหรับการทำนายคุณภาพดิน (Deep Learning Models for Soil Quality Prediction) มีความสำคัญอย่างยิ่งในระบบเกษตรอัจฉริยะ โดยมีรายละเอียดสาระสำคัญดังนี้

ระบบ SoilAI ใช้สถาปัตยกรรมโครงข่ายประสาทเทียมแบบคอนโวลูชัน (Convolutional Neural Network; CNN) ร่วมกับอัลกอริทึมการถดถอย (Regression Algorithms):
- สถาปัตยกรรม Backbone: ใช้โครงข่ายน้ำหนักเบาประสิทธิภาพสูง เช่น MobileNetV3 หรือ EfficientNet-B0 ที่สามารถประมวลผลบนสมาร์ทโฟนได้โดยตรง (On-device Edge Inference)
- การสกัดคุณลักษณะ (Feature Extraction): ดึงลักษณะพื้นผิวของเม็ดดิน (Soil texture), การรวมตัวของก้อนดิน (Aggregate structure), และการกระจายตัวของโทนสี
- ชั้นทำนายค่าต่อเนื่อง (Multi-output Regression Head): ทำนายค่าพารามิเตอร์คุณภาพดินพร้อมกัน ได้แก่ ค่าความชื้นในดิน (Moisture \%), ค่าความเป็นกรด-ด่าง (pH), และค่าปริมาณอินทรียวัตถุ (OM \%)
- การประเมินความแม่นยำของโมเดล: วัดด้วยค่าสัมประสิทธิ์การตัดสินใจ ($R^2 > 0.85$) และค่ารากที่สองของความคลาดเคลื่อนเฉลี่ย (RMSE ต่ำ) เทียบกับผลแล็บจริง

\begin{techbox}[ข้อกำหนดทางเทคนิคสถาปัตยกรรม AIoT]
การเชื่อมต่อระบบเซนเซอร์ไร้สายต้องรักษาความมั่นคงปลอดภัยของข้อมูล มีการเข้ารหัสระดับ AES-128 และบริหารจัดการพลังงานให้โหนดเซนเซอร์ทำงานได้อย่างต่อเนื่องไม่น้อยกว่า 3 ปี
\end{techbox}

\section{ขั้นตอนและวิธีปฏิบัติการทดลอง}
\index{ขั้นตอนการทดลองประจำบทที่ 6}

\subsection{ขั้นตอนการถ่ายภาพดินด้วยสมาร์ทโฟนตามระเบียบวิธีมาตรฐาน SoilAI}
\begin{sensorbox}[ขั้นตอนการถ่ายภาพดินด้วยสมาร์ทโฟนตามระเบียบวิธีมาตรฐาน SoilAI]
1. เกลี่ยผิวดินบริเวณที่จะถ่ายให้เรียบสม่ำเสมอในพื้นที่สี่เหลี่ยมขนาด 15x15 เซนติเมตร
2. วางบัตรสีอ้างอิงมาตรฐาน (Calibration Card) บนผิวดินบริเวณขอบภาพ
3. ถือสมาร์ทโฟนขนานกับผิวดินในแนวดิ่ง (Top-down view) ที่ระยะห่างคงที่ 20-25 เซนติเมตร
4. ปิดแฟลชกล้องและหลีกเลี่ยงการทอดเงาของผู้ถ่ายลงบนผิวดิน
5. บันทึกภาพในโหมดความละเอียดสูง บันทึกพิกัด GPS อัตโนมัติ
\end{sensorbox}

\subsection{การรันโมเดลทำนายคุณภาพดินบนเว็บแอปพลิเคชัน SoilAI}
\begin{sensorbox}[การรันโมเดลทำนายคุณภาพดินบนเว็บแอปพลิเคชัน SoilAI]
1. เปิดเว็บแอปพลิเคชัน Soil Analysis AI ผ่านเบราว์เซอร์สมาร์ทโฟน
2. อัปโหลดไฟล์ภาพถ่ายดินเข้าสู่ระบบ
3. อัลกอริทึมตรวจจับบัตรสีอัตโนมัติและปรับเทียบสี (Auto Color Normalization)
4. โมเดล Deep Learning ประมวลผลภาพภายในเวลาไม่เกิน 3 วินาที
5. หน้าจอแสดงผลการวินิจฉัย: ค่าความชื้นในดิน, ระดับความเป็นกรด-ด่าง, และระดับปริมาณอินทรียวัตถุ พร้อมคำแนะนำการปรับปรุงดิน
\end{sensorbox}

\subsection{การตรวจสอบความถูกต้องของโมเดลด้วยชุดข้อมูลทดสอบ}
\begin{sensorbox}[การตรวจสอบความถูกต้องของโมเดลด้วยชุดข้อมูลทดสอบ]
1. รวบรวมภาพถ่ายดินจำนวน 100 ตัวอย่างที่ทราบผลวิเคราะห์จากแล็บปฐพีวิทยา
2. ป้อนภาพเข้าสู่โมเดลเพื่อคำนวณค่าทำนาย ($y_{\text{pred}}$)
3. เปรียบเทียบกับค่าจริงจากห้องแล็บ ($y_{\text{true}}$)
4. พล็อตเส้นตรงความสัมพันธ์ (Scatter plot) และคำนวณค่า $R^2$ และ Mean Absolute Error (MAE)
5. หากค่า MAE ของ pH ต่ำกว่า 0.35 หน่วย ถือว่าโมเดลมีความแม่นยำยอมรับได้ในระดับการจัดการแปลง
\end{sensorbox}

\section{ตารางบันทึกผลการทดลองและการอภิปรายผล}
\begin{table}[H]
\centering\small
\caption{ตารางบันทึกผลการตรวจวัดและข้อมูลพารามิเตอร์ประจำบทที่ 6}
\label{tab:ch06_datasheet}
\begin{tabularx}{\textwidth}{p{0.3\textwidth} p{0.3\textwidth} X}
\toprule
\textbf{พารามิเตอร์ / การทดสอบ} & \textbf{ค่าที่ตรวจวัดได้} & \textbf{การแปลผลและการวิเคราะห์} \\
\midrule
จุดตรวจวัดที่ 1 (บริเวณดินดอน) & ค่าความชื้นอยู่ในเกณฑ์ปกติ & ปริมาณน้ำเพียงพอต่อการเจริญ \\
จุดตรวจวัดที่ 2 (บริเวณที่ลุ่ม) & มีการสะสมความชื้นสูง & ควรชะลอการให้น้ำเพื่อป้องกันรากเน่า \\
สถานะสัญญาณโครงข่าย & RSSI: -85 dBm, SNR: +8 dB & คุณภาพสัญญาณ LoRaWAN ยอดเยี่ยม \\
\bottomrule
\end{tabularx}
\end{table}

\section{คำถามทบทวนและแบบฝึกหัดท้ายบท}
\begin{enumerate}
    \item เหตุใดการปรับเทียบสีด้วยแผ่นเทียบสี (Color Calibration Card) จึงมีความจำเป็นอย่างยิ่งในการวิเคราะห์ภาพถ่ายดิน
    \item จงอธิบายว่าค่าความสว่าง $L^*$ ในระบบสี CIE L*a*b* มีความสัมพันธ์เชิงตัวเลขกับปริมาณความชื้นและอินทรียวัตถุในดินอย่างไร
    \item หากสภาพแสงขณะถ่ายภาพมีเงาตกทอดครึ่งหนึ่งของภาพ จะส่งผลต่อความแม่นยำของโมเดล Deep Learning อย่างไร และมีวิธีแก้ไขในขั้นตอน Pre-processing อย่างไร
\end{enumerate}
